\chapter{Preface}

This book began with dissatisfaction over explanations that had become familiar enough to feel obvious: that the one to whom God gave the disclosure is Jesus Christ, that the servants are Christians generally, that \textit{en tachei} means soon, and that the opening lines describe a simple linear chain of communication.

The question pursued here is not whether those readings are possible. Many of them are possible, and some are nearly convincing. The question is whether they are governed by an interpretation system strong enough to make the choices predictable, consistent, economical, and capable of excluding alternatives.

The argument of this book is that Revelation 1:1--3 is more compressed than most translations allow readers to see. Its reduced form has encouraged interpreters to supply missing constituents by plausibility, intertextual appeal, theological convenience, or inherited habit. But the passage also contains an expanded form in which ellipsis, verbal structure, and arrangement constrain what must be supplied.

This is a study of those constraints.
